\section{Technologies de collecte de données en agriculture de précision}

\subsection{Les capteurs de sol et de climat}
En agriculture de précision, les capteurs constituent des outils essentiels car ils permettent de mesurer les paramètres du sol et de l'environnement. Les capteurs de sol mesurent notamment la température, l'humidité, la EC, le pH ainsi que les teneurs en N, P, K. 
Les appareils combinés, comme les capteur 8-en-1 regroupent plusieurs fonctions dans un même dispositif facilitant la collecte de données.

Les données météorologiques telles que l'humidité de l'air, la température et la vitesse du vent permettent d'estimer les besoins en eau des plantes. 
L'association des données du climat et du sol contribue ainsi à une meilleure gestion de l'irrigation et de la fertilisation.

\begin{figure}[H]
    \centering
    \imgstd{images/01_capteur_8_en_1.jpg}
    \caption{Capteur de sol 8-en-1. }
    \label{fig:8_en_1}
\end{figure}

\subsection{Les drones et l’imagerie aérienne en agriculture}
Les drones équipés de capteurs multispectraux ou optiques permettent de surveiller l'état des cultures et de détecter les zones de stress nutritif et/ou hydrique. Toutefois, leur utilisation nécessite des compétences techniques et des équipements spécifiques. Cependant cet outil n'a pas été utilisée dans notre expérimentation.

\subsection{Les donnees météorologiques et les API météo}
Les données climatiques telles que, la $T_{max}$ et $T_{min}$, la $RH_{max}$ et $RH_{min}$ ainsi que la vitesse du vent sont indispensables à l'estimation des besoins en eau des cultures. 
Pendant l'expérimentation, ces données ont été relevées à partir de l'application AccuWeather du smartphone. A la fin de l'expérimentation, l'API OpenWeather a été retenue pour automatiser la récupération des paramètres météorologiques de l'application conçue.

\subsection{Limites et défis de fiabilité des capteurs low-cost}
Malgré leur accessibilité, les \textbf{capteurs à faible coût} peuvent présenter des erreurs de mesure. La calibration, l'entretien et la vérification régulière des données sont donc nécessaires afin d'assurer leur fiabilité et d'améliorer la précision des décisions agricoles.

